\documentclass[a4paper]{article}

\usepackage[margin=7mm]{geometry}
\usepackage{array}
\usepackage[table]{xcolor}
\usepackage{listings}
\usepackage{tabularx}

\pagestyle{empty}
\setlength{\parindent}{0pt}
\setlength{\tabcolsep}{2.5pt}
\renewcommand{\arraystretch}{1.03}

\definecolor{cabecera}{HTML}{D9EAF7}
\definecolor{borde}{HTML}{5B6573}

\lstdefinelanguage{OracleSQL}{
  language=SQL,
  morekeywords={CONCAT,SUBSTR,UPPER},
  sensitive=false
}

\lstset{
  language=OracleSQL,
  basicstyle=\ttfamily\fontsize{7pt}{8pt}\selectfont,
  keywordstyle=\bfseries,
  stringstyle=\itshape,
  breaklines=true,
  breakatwhitespace=true,
  columns=fullflexible,
  keepspaces=true,
  showstringspaces=false,
  aboveskip=1pt,
  belowskip=0pt,
  frame=none
}

\newcommand{\consultas}{pauta-es1-consultas-sql.sql}
\newcommand{\codigo}[2]{%
  \begin{minipage}[t]{\linewidth}
    \vspace{0.5pt}
    \lstinputlisting[firstline=#1,lastline=#2]{\consultas}
    \vspace{0.5pt}
  \end{minipage}%
}

\begin{document}
\color{black}
\arrayrulecolor{borde}

\begin{minipage}[t]{\textwidth}
\begin{tabularx}{\linewidth}{|>{\centering\arraybackslash\bfseries\fontsize{7.5pt}{8pt}\selectfont}m{7mm}|X|}
\hline
\rowcolor{cabecera}
Req. & \centering\arraybackslash\bfseries\fontsize{7.5pt}{8pt}\selectfont Código SQL \\
\hline
1 & \codigo{9}{16} \\
\hline
2 & \codigo{22}{28} \\
\hline
3 & \codigo{34}{42} \\
\hline
\end{tabularx}
\end{minipage}

\vspace{3mm}

\begin{minipage}[t]{\textwidth}
\begin{tabularx}{\linewidth}{|>{\centering\arraybackslash\bfseries\fontsize{7.5pt}{8pt}\selectfont}m{7mm}|X|}
\hline
\rowcolor{cabecera}
Req. & \centering\arraybackslash\bfseries\fontsize{7.5pt}{8pt}\selectfont Código SQL \\
\hline
4 & \codigo{48}{52} \\
\hline
5 & \codigo{58}{68} \\
\hline
6 & \codigo{75}{92} \\
\hline
\end{tabularx}
\end{minipage}

\clearpage

\begin{minipage}[t]{\textwidth}
\begin{tabularx}{\linewidth}{|>{\centering\arraybackslash\bfseries\fontsize{7.5pt}{8pt}\selectfont}m{7mm}|X|}
\hline
\rowcolor{cabecera}
Req. & \centering\arraybackslash\bfseries\fontsize{7.5pt}{8pt}\selectfont Código SQL \\
\hline
7 & \codigo{99}{113} \\
\hline
\end{tabularx}
\end{minipage}

\vspace{4mm}

{\fontsize{6.6pt}{7.5pt}\selectfont
\begin{tabularx}{\textwidth}{|>{\centering\arraybackslash\ttfamily\bfseries}m{13mm}|X|>{\centering\arraybackslash\ttfamily\bfseries}m{13mm}|X|}
\hline
\rowcolor{cabecera}
\multicolumn{4}{|c|}{\bfseries Formatos principales de fecha para \texttt{TO\_CHAR(SYSDATE, formato)}} \\
\hline
DD & Día del mes en número: 01--31 & D & Día de la semana en número: 1--7 \\
\hline
DY & Nombre abreviado del día: LUN & DAY & Nombre completo del día: LUNES \\
\hline
DDD & Día del año: 001--366 & MM & Mes en número: 01--12 \\
\hline
MON & Nombre abreviado del mes: SEP & MONTH & Nombre completo del mes: SEPTIEMBRE \\
\hline
RM & Mes en números romanos: I--XII & Q & Trimestre del año: 1--4 \\
\hline
W & Semana del mes: 1--5 & WW & Semana del año: 01--53 \\
\hline
IW & Semana ISO del año: 01--53 & YYYY & Año con cuatro dígitos: 2026 \\
\hline
YY & Últimos dos dígitos del año: 26 & RR & Año de dos dígitos con ajuste de siglo \\
\hline
RRRR & Año de dos o cuatro dígitos & CC & Número del siglo \\
\hline
HH12 & Hora en formato de 12 horas: 01--12 & HH24 & Hora en formato de 24 horas: 00--23 \\
\hline
MI & Minutos: 00--59 & SS & Segundos: 00--59 \\
\hline
SSSSS & Segundos transcurridos desde medianoche & AM/PM & Indicador antes o después del mediodía \\
\hline
J & Número de día juliano & FM & Elimina ceros y espacios de relleno \\
\hline
\textnormal{"texto"} & Inserta texto literal dentro de la fecha & / : - , & Separadores que se muestran en el resultado \\
\hline
\rowcolor{cabecera}
\multicolumn{4}{|c|}{\bfseries Formatos principales de números para \texttt{TO\_CHAR(numero, formato)}} \\
\hline
9 & Dígito opcional; deja un espacio si falta & 0 & Dígito obligatorio; completa con cero \\
\hline
G & Separador de miles definido por NLS & D & Separador decimal definido por NLS \\
\hline
\$ & Símbolo de dólar & L & Símbolo monetario local definido por NLS \\
\hline
S & Signo positivo o negativo & MI & Signo negativo colocado al final \\
\hline
PR & Negativos entre ángulos & FM & Elimina espacios y ceros de relleno \\
\hline
EEEE & Notación científica & V & Multiplica por una potencia de diez \\
\hline
\end{tabularx}
}

\end{document}
